% Compilar con XeLaTeX (dos pasadas) desde este directorio:
% xelatex -interaction=nonstopmode -halt-on-error pasteles_rsi.tex
% No requiere imágenes externas ni acceso a red.
\documentclass[11pt,a4paper]{article}
\usepackage[margin=21mm,top=20mm,bottom=20mm,headheight=15pt]{geometry}
\usepackage{fontspec}
\setmainfont{TeX Gyre Pagella}
\setsansfont{TeX Gyre Heros}
\setmonofont[Scale=0.85]{DejaVu Sans Mono}
\usepackage[spanish,es-noshorthands,es-nodecimaldot]{babel}
\usepackage{microtype}
\usepackage{amsmath,amssymb}
\usepackage{booktabs,tabularx,array}
\usepackage{enumitem}
\usepackage{xcolor}
\usepackage{tikz}
\usetikzlibrary{arrows.meta,positioning,fit,calc}
\usepackage[most]{tcolorbox}
\usepackage{fancyhdr}
\usepackage{titlesec}
\usepackage{hyperref}
\definecolor{ink}{HTML}{173242}
\definecolor{teal}{HTML}{087F8C}
\definecolor{muted}{HTML}{526674}
\definecolor{paper}{HTML}{F0F5F6}
\definecolor{amber}{HTML}{9B651B}
\definecolor{amberbg}{HTML}{FCF6EB}
\color{ink}
\hypersetup{colorlinks=true,linkcolor=teal,urlcolor=teal,
  pdftitle={Pasteles: agentes personales y mejora recursiva},
  pdfauthor={Proyecto Pasteles},pdfsubject={Arquitectura, ontología, RSI y ejemplos evaluados de conversación}}
\pagestyle{fancy}
\fancyhf{}
\fancyhead[L]{\sffamily\small\bfseries PASTELES\quad\textcolor{teal}{/}\quad INFORME DEL PROYECTO}
\fancyhead[R]{\sffamily\small\textcolor{muted}{RSI · 2026}}
\fancyfoot[L]{\sffamily\footnotesize\textcolor{muted}{Demostración privada · catálogo y pagos simulados}}
\fancyfoot[R]{\sffamily\footnotesize\thepage}
\renewcommand{\headrulewidth}{0.4pt}
\renewcommand{\footrulewidth}{0pt}
\titleformat{\section}{\sffamily\Large\bfseries\color{ink}}{\textcolor{teal}{\thesection}}{0.65em}{}
\titleformat{\subsection}{\sffamily\normalsize\bfseries\color{teal}}{}{0pt}{}
\titlespacing*{\section}{0pt}{0pt}{12pt}
\titlespacing*{\subsection}{0pt}{13pt}{5pt}
\setlength{\parindent}{0pt}
\setlength{\parskip}{6pt}
\setlength{\emergencystretch}{2em}
\setlist[itemize]{leftmargin=1.3em,itemsep=3pt,topsep=3pt}
\setlist[enumerate]{leftmargin=1.55em,itemsep=5pt,topsep=3pt}
\newcolumntype{Y}{>{\raggedright\arraybackslash}X}
\newcolumntype{P}[1]{>{\raggedright\arraybackslash}p{#1}}
\newcommand{\code}[1]{\texttt{\detokenize{#1}}}
\newcommand{\eyebrow}[1]{{\sffamily\small\bfseries\color{teal}#1\par}}
\newcommand{\source}[1]{\par{\footnotesize\color{muted}Fuente: #1.\par}}
\newtcolorbox{nota}[1][]{colback=paper,colframe=teal,boxrule=0pt,
  borderline west={2pt}{0pt}{teal},arc=0pt,left=10pt,right=10pt,top=8pt,bottom=8pt,#1}
\newtcolorbox{limite}[1][]{colback=amberbg,colframe=amber,boxrule=0pt,
  borderline west={2pt}{0pt}{amber},arc=0pt,left=10pt,right=10pt,top=8pt,bottom=8pt,#1}
\newtcolorbox{dialogo}[2][]{enhanced,colback=paper,colframe=teal,
  colbacktitle=teal,coltitle=white,fonttitle=\sffamily\small\bfseries,
  fontupper=\small,boxrule=0pt,arc=0pt,left=9pt,right=9pt,top=8pt,bottom=8pt,
  title={#2},#1}
\newcommand{\habla}[2]{\textbf{#1}\enspace #2\par\smallskip}
\tikzset{caja/.style={draw=teal!50,fill=paper,rounded corners=2pt,
  align=center,font=\sffamily\small,inner sep=7pt,minimum height=12mm},
  flecha/.style={-{Stealth[length=2mm]},thick,draw=teal},
  rotulo/.style={font=\sffamily\scriptsize,fill=white,inner sep=2pt,align=center}}

\begin{document}
\thispagestyle{empty}
\eyebrow{PROYECTO PARA EL HACKATÓN DE RECURSIVE SELF-IMPROVEMENT}
\vspace{12mm}
{\sffamily\fontsize{44}{48}\selectfont\bfseries PASTELES\par}
\vspace{4mm}
{\sffamily\fontsize{23}{29}\selectfont Agentes personales y\par mejora recursiva verificable\par}
\vspace{6mm}
{\large Un vendedor conversacional que evoluciona su harness\par y las skills que construyen a su sucesor.\par}
\vspace{7mm}
\textcolor{teal}{\rule{\linewidth}{1.5pt}}\par
\vspace{5mm}
\textbf{Resumen.} Pasteles es una demostración privada de venta de pasteles por
Telegram. Un agente personal conserva preferencias del cliente; un agente
vendedor recibe únicamente los datos autorizados para la conversación. Un modelo
pequeño interpreta solicitudes, mientras el código calcula precios, verifica
condiciones y controla pagos simulados.

El ciclo RSI evalúa una versión, propone un cambio acotado a su harness y actualiza
la skill del constructor. Una compuerta independiente decide si aceptar el
sucesor. La versión aceptada construye la siguiente generación, con linaje,
registro de rechazos y posibilidad de volver atrás.
Este informe añade cuatro ejemplos de conversación y la medición de una
corrección del diálogo motivada por una captura real del equipo.
\vspace{4mm}
\begin{nota}
\begin{tabularx}{\linewidth}{@{}YYY@{}}
{\sffamily\Large\bfseries 495} & {\sffamily\Large\bfseries 6/15 $\to$ 13/15} & {\sffamily\Large\bfseries 7/15 $\to$ 15/15}\\[3pt]
Pruebas de software aprobadas & RSI: primer sucesor con Claude real & Experiencia: nueva evaluación con Claude real\\
\end{tabularx}
\end{nota}
\begin{limite}
\textbf{Alcance de la evidencia.} Son instrumentos distintos sobre guiones
sintéticos. El catálogo y los pagos son ficticios. La reserva histórica alcanzó
38/45 objetivos y corresponde al código anterior; la corrección reciente aún no
tiene una reserva independiente. No se presentan propuestas como ventas.
\end{limite}
\vfill
{\sffamily\small\textbf{Corte documental:} 9 de noviembre de 2026, 18:07 de Monterrey\par
\textbf{Código de referencia:} \code{a242a35}, fases 5 a 9 registradas\par
\textbf{Gasto acumulado al corte:} USD 0.045687 de un máximo de USD 5\par
\textbf{Guía de lectura:} arquitectura y RSI en las secciones 1--4;
ejemplos y nueva medición en 5--7; negocio, seguridad, panel y operación en 8--12.}

\newpage
\section{Del mensaje a una cotización verificable}
El problema inicial fue concreto: el modelo podía omitir una nueva cotización
después de un cambio de fecha o presupuesto. Reforzar el prompt mejoró parte de la
conducta, pero no garantizó que se invocara la herramienta. El nuevo harness lleva
esa decisión a una máquina de estado explícita.

\begin{center}
\begin{tikzpicture}[node distance=7mm]
\node[caja,text width=29mm] (u) {Cliente\\Telegram};
\node[caja,text width=33mm,right=of u] (a) {Adaptador\\autenticado};
\node[caja,text width=33mm,right=of a] (h) {Harness\\de conversación};
\node[caja,text width=29mm,right=of h] (b) {Motor\\de negocio};
\node[caja,text width=40mm,below=13mm of h] (p) {Proxy local\\presupuesto persistente};
\node[caja,text width=40mm,below=13mm of a,xshift=-15mm] (m) {Modelo pequeño\\extracción JSON};
\draw[flecha] (u)--(a); \draw[flecha] (a)--(h); \draw[flecha] (h)--(b);
\draw[flecha] (h)--node[rotulo,right]{reserva y\\validación}(p);
\draw[flecha] (p)--node[rotulo,below]{API}(m);
\end{tikzpicture}
\end{center}
{\small\textcolor{muted}{Esquema lógico: toda inferencia del modelo pasa por el proxy y la frontera LLM.}}

\subsection{Responsabilidades separadas}
\begin{tabularx}{\linewidth}{@{}P{32mm}Y@{}}
\toprule
\textbf{Componente} & \textbf{Qué hace}\\
\midrule
Adaptador & Autentica usuario y chat, serializa mensajes, consume comandos y evita que un error se delegue al agente general.\\
Modelo & Extrae intención y campos permitidos: sabor, asistentes, fecha, entrega, restricciones y presupuesto.\\
Harness & Conserva pedido y preguntas pendientes, decide qué falta, invalida la cotización anterior y ejecuta la política versionada.\\
Negocio y pagos & Calculan importes en centavos, aplican disponibilidad y restricciones, y registran decisiones humanas de simulación.\\
\bottomrule
\end{tabularx}

\subsection{Un turno ordinario}
\begin{enumerate}
\item Se invalida el pago anterior antes de esperar una respuesta del modelo.
\item Se solicita una extracción JSON, con hasta 768 tokens de salida. El extractor
recibe pedido, preguntas pendientes y última respuesta acotada; una entrada
inválida borra el estado ambiguo y obliga a recoger los datos de nuevo.
\item El harness combina los datos del turno con las preferencias autorizadas que
permita la versión; los datos actuales tienen prioridad.
\item Se comprueba la disponibilidad y el motor produce una cotización nueva cuando
hay datos suficientes. El texto final y sus importes se
construyen a partir de ese resultado; el modelo no aprueba el pago.
\end{enumerate}

\begin{nota}
\textbf{Ejemplo.} Chocolate para 10 personas, entrega en centro: pastel de
\$650 MXN y entrega de \$60 MXN, total de \$710 MXN. Si el presupuesto cambia a
\$700, el motor devuelve \code{excede_presupuesto} y deshabilita el pago.
\end{nota}
\source{\code{pasteles/arnes.py}, \code{pasteles/negocio.py} y \code{pasteles/telegram_rsi.py}}

\newpage
\section{Ontología: personas, agentes y conversaciones}
La ontología es el vocabulario compartido que permite identificar quién es dueño
de un dato, qué agente puede actuar y qué significa cada mensaje. Se implementa
como contratos tipados y un enrutador; los permisos se establecen fuera del texto
que genera el modelo.

\begin{center}
\begin{tikzpicture}[node distance=9mm]
\node[caja,text width=31mm] (pc) {Persona cliente};
\node[caja,text width=31mm,below=of pc] (ac) {Agente personal};
\node[caja,text width=31mm,below=of ac] (ctx) {Contexto privado};
\node[caja,text width=38mm,right=17mm of ac] (conv) {Conversación\\para cotizar};
\node[caja,text width=31mm,right=17mm of conv] (av) {Agente vendedor};
\node[caja,text width=31mm,above=of av] (pv) {Persona pastelería};
\node[caja,text width=31mm,below=of av] (sk) {Skills versionadas};
\draw[flecha] (pc)--(ac); \draw[flecha] (ac)--(ctx);
\draw[flecha] (pv)--(av); \draw[flecha] (av)--(sk);
\draw[flecha] (ac)--(conv); \draw[flecha] (av)--(conv);
\draw[flecha,dashed] (ctx.east)--node[rotulo,below]{solo campos\\autorizados}(conv.south);
\end{tikzpicture}
\end{center}

\begin{tabularx}{\linewidth}{@{}P{35mm}Y@{}}
\toprule
\textbf{Entidad} & \textbf{Contrato}\\
\midrule
Persona y agente & Identidad opaca, propietario, rol, capacidades y referencias a skills. La autenticación procede del adaptador.\\
Contexto & Preferencias estructuradas de un cliente. Guardarlas y compartirlas son permisos distintos.\\
Conversación & Participantes y finalidad explícita: \code{cotizar_pastel}.\\
Autorización & Concesión para campos concretos, origen, destino y conversación determinados.\\
Sobre & Mensaje con origen, destino, tipo, datos y correlación. No crea permisos ni reemplaza una skill.\\
\bottomrule
\end{tabularx}

\subsection{Tres mensajes con significado preciso}
\code{pedido} comunica la solicitud del agente personal al vendedor.
\code{cotizacion} devuelve solicitud y resultado vinculados por correlación; el
enrutador comprueba que coincidan con el motor de negocio.
\code{preferencias_compartidas} transporta únicamente el subconjunto concedido.

\subsection{Memoria para un cliente frecuente}
Se admiten sabor, modalidad, zona e idioma mediante valores enumerados. No se
guardan instrucciones libres, precios, fechas ni restricciones alimentarias en
esta memoria. Los archivos tienen nombres opacos derivados con HMAC, contenido
firmado y permisos privados; el contexto de un cliente no se entrega al constructor.

\begin{nota}
\code{/recordar sabor=chocolate modalidad=recoger} autoriza guardar.\\
\code{/compartir sabor modalidad} autoriza compartir en esa conversación.\\
\code{/new} revoca la concesión; \code{/olvidar} borra las preferencias propias.
\end{nota}
\source{\code{pasteles/ontologia.py}, \code{pasteles/contexto_clientes.py} y \code{docs/ONTOLOGIA.md}}

\newpage
\section{Cómo funciona la mejora recursiva}
La unidad que evoluciona es una versión del sistema de conversación. El modelo
base permanece fijo. Cada versión incluye una política ejecutable del harness,
una skill de vendedor y una skill de constructor:
\[
 V_n=(H_n,\,S_n^{\mathrm{vendedor}},\,S_n^{\mathrm{constructor}}).
\]
El constructor utiliza la versión aceptada y resultados agregados de desarrollo
para producir un candidato. Su propia skill actualizada pasa al sucesor; por eso
la siguiente generación no parte otra vez del constructor inicial.

\begin{center}
\begin{tikzpicture}[node distance=8mm]
\node[caja,text width=39mm] (v) {Versión aceptada $V_n$};
\node[caja,text width=39mm,right=of v] (e) {Evaluación y\\fallos agregados};
\node[caja,text width=39mm,right=of e] (c) {Constructor con\\skill heredada};
\node[caja,text width=39mm,below=12mm of c] (p) {Candidato $V_{n+1}$\\una hipótesis};
\node[caja,text width=39mm,left=of p] (j) {Jueces\\independientes};
\node[caja,text width=39mm,left=of j] (a) {Aceptar o conservar\\la versión anterior};
\draw[flecha] (v)--(e); \draw[flecha] (e)--(c);
\draw[flecha] (c)--(p); \draw[flecha] (p)--(j); \draw[flecha] (j)--(a);
\draw[flecha] (a)--node[rotulo,left]{herencia}(v);
\end{tikzpicture}
\end{center}

\subsection{Un espacio de cambios acotado}
Un candidato modifica exactamente una decisión: cotización automática, uso de
preferencias compartidas, una o dos preguntas por turno, cierre consultivo o
estrategia comercial. Debe actualizar la skill del constructor; también puede
revisar la del vendedor. La salida es un parche JSON que el orquestador ensambla y
valida. No se ejecuta código generado. Si el JSON es inválido, se permite una sola
corrección de formato, con el mismo presupuesto.

\subsection{Compuerta de aceptación}
La decisión compara comprobaciones pareadas bajo el mismo instrumento:
\[
 A=E\land O\land(S=1)\land(R=0)\land(M>0).
\]
Aquí $E$ significa evaluación íntegra y completa; $O$, costo y latencia dentro de
los criterios; $S$, proporción de controles de seguridad aprobados; $R$, número de
regresiones; y $M$, mejoras funcionales o comerciales. Una mejora comercial puede
acompañarse de empate funcional, pero no compensar una regresión.

\begin{limite}
\textbf{Límite de autoridad.} El candidato no puede modificar juez, corpus,
umbrales, presupuesto, modelo, permisos ni herramientas. Un empate detiene la
ronda. Los resultados negativos se conservan y la versión anterior sigue activa.
\end{limite}
El manifiesto vincula contenido y padre mediante huellas. La bitácora identifica
propuestas, decisiones y publicaciones; el rollback admite versiones iniciales o
previamente aceptadas.
\source{\code{pasteles/constructor_rsi.py}, \code{pasteles/orquestador_rsi.py} y \code{pasteles/versiones_rsi.py}}

\newpage
\section{Evaluaciones y evidencia observada}
Los instrumentos separan comportamiento del harness, uso del modelo y operación
del canal. Desarrollo sirve para proponer sucesores; la reserva observa el
comportamiento en casos apartados y no se transmite al constructor. Los hashes fijan corpus,
criterios y jueces; revisiones diferentes no se comparan como el mismo experimento.

\begin{tabularx}{\linewidth}{@{}P{38mm}P{44mm}Y@{}}
\toprule
\textbf{Evidencia} & \textbf{Resultado} & \textbf{Qué permite afirmar}\\
\midrule
Suite de software actual & 495 pruebas aprobadas & Contratos y regresiones cubiertos por pruebas; no son 495 conversaciones reales.\\
RSI determinista & 13/18 $\to$ 18/18; seguridad 119/119 & Mejora de política con reglas y extracciones preparadas.\\
Sucesor real & 6/15 $\to$ 13/15; seguridad 119/119 & Siete mejoras pareadas, sin regresiones, con propuesta y extracción mediante Claude.\\
Reserva RSI histórica & 36/36 conversaciones; 38/45 objetivos; seguridad 357/357 & Casos apartados del desarrollo, con siete objetivos incumplidos; evaluó el código anterior.\\
Comercial sintética & Cierre 0/5 $\to$ 5/5; valor 0/2 $\to$ 2/2; seguridad 61/61 & Preparación del cierre y aprovechamiento de oportunidades definidas.\\
\bottomrule
\end{tabularx}
\source{\code{docs/VERIFICACION.md} y reportes agregados indicados en la sección \ref{sec:fuentes}}

\subsection{El historial incluye interrupciones y rechazos}
\begin{enumerate}
\item El ciclo determinista aceptó cuatro generaciones y rechazó la quinta por no
mejorar. Su constructor fue de reglas: demuestra el mecanismo del ciclo.
\item El primer ciclo real publicó un sucesor, pero terminó interrumpido al fallar
el formato JSON de una propuesta posterior.
\item El segundo intento real se interrumpió antes de publicar. Después de reducir
la propuesta a un parche pequeño, un tercer ciclo evaluó y rechazó un empate.
\end{enumerate}
La versión real aceptada habilita recotización automática. Las políticas de
preferencias automáticas, cierre consultivo y mayor valor todavía no fueron
promovidas por el ciclo real.

\subsection{Costo, latencia y lectura de los denominadores}
La reserva consumió \textbf{USD 0.0058074}, con latencia p95 de \textbf{3.79 segundos}
por caso. Al cerrar esa etapa, el gasto fue \textbf{USD 0.0413926 de 5}, sin reservas
pendientes. Tras la corrección de experiencia y el uso intermedio del bot, el
acumulado llegó a \textbf{USD 0.045687}; se detalla la nueva medición en la sección 7.

\begin{limite}
Los 18 objetivos deterministas y los 15 objetivos reales tienen distinto alcance:
tres oportunidades corresponden a entradas estructuradas inválidas. No deben
compararse como muestras equivalentes. Tampoco los 495 tests ni el 100\% de los
controles cubiertos demuestran superioridad estadística o inmunidad universal.
\end{limite}

\newpage
\section{Ejemplos: intención y respuestas breves}
Una captura de Telegram mostró fallos que los benchmarks previos no explicaban
bien: un «no» alimentario cancelaba la compra y una petición informal recibía una
respuesta de seguridad. Se reprodujeron en un instrumento nuevo.

\begin{nota}
\textbf{Cómo leer estos ejemplos.} La captura motivó el cambio. Los diálogos de
estas dos páginas son \textbf{extractos de reproducciones sintéticas con Claude
real}, antes y después de la corrección; no son una continuación del chat privado.
Se omiten el prefijo repetido de simulación y detalles sin relación con el fallo.
\end{nota}

\subsection{1. «No» responde a una pregunta; no cancela automáticamente}
\textbf{Situación:} chocolate para cinco personas, jueves, entrega en centro.
El único dato pendiente es si existe una restricción alimentaria.

\begin{tcbraster}[raster columns=2,raster column skip=4mm,raster equal height=rows]
\begin{dialogo}[colback=amberbg,colbacktitle=amber]{ANTES · CLAUDE REAL}
\habla{Bot}{¿Hay alguna restricción alimentaria?}
\habla{Cliente}{no}
\habla{Bot}{De acuerdo. Gracias por visitarnos.}
\textcolor{amber}{\textbf{Resultado:} se pierde el pedido.}
\end{dialogo}
\begin{dialogo}{DESPUÉS · CLAUDE REAL}
\habla{Bot}{¿Hay alguna restricción alimentaria?}
\habla{Cliente}{no}
\habla{Bot}{Chocolate intenso · 12 porciones para 5 personas. Total: \$710.00 MXN.}
\textcolor{teal}{\textbf{Resultado:} conserva el pedido y cotiza.}
\end{dialogo}
\end{tcbraster}

\textbf{Qué cambió.} El harness recuerda la pregunta pendiente y pregunta por
restricciones por separado. Una negación mínima ante esa única pregunta significa
«ninguna». «No gracias, cancela el pedido» sigue siendo una cancelación.

\subsection{2. Los errores de escritura no son una instrucción maliciosa}
\begin{tcbraster}[raster columns=2,raster column skip=4mm,raster equal height=rows]
\begin{dialogo}[colback=amberbg,colbacktitle=amber]{ANTES · CLAUDE REAL}
\habla{Cliente}{tengo ambre}
\habla{Bot}{Puedo ayudarte con el catálogo y una cotización. No puedo cambiar reglas ni acceder a datos de otras personas.}
\end{dialogo}
\begin{dialogo}{DESPUÉS · CLAUDE REAL}
\habla{Cliente}{tengo ambre}
\habla{Bot}{¿Qué sabor prefieres? ¿Para cuántas personas?}
\end{dialogo}
\end{tcbraster}

\textbf{Qué cambió.} El contrato de extracción distingue una expresión informal
de interés de un intento de cambiar permisos, precios o reglas. La variante
«vendene un pastel» también recibió preguntas comerciales en la medición corregida.
\source{Experiencia v1, casos \code{continuidad_no} y \code{coloquial_hambre}; reportes reales de la sección \ref{sec:fuentes}}

\newpage
\section{Ejemplos: disponibilidad y orientación de pago}
La conversación debe resolver primero el impedimento relevante. Pedir todos los
datos de un producto agotado o rechazar una pregunta de pago añade fricción sin
ayudar al cliente a decidir.

\subsection{3. Comprobar el catálogo antes de seguir preguntando}
\begin{tcbraster}[raster columns=2,raster column skip=4mm,raster equal height=rows]
\begin{dialogo}[colback=amberbg,colbacktitle=amber]{ANTES · CLAUDE REAL}
\habla{Cliente}{quiero un pastel de fresa}
\habla{Bot}{¿Para cuántas personas? ¿Para qué fecha lo necesitas?}
\textcolor{amber}{\textbf{Resultado:} solicita más datos sin advertir todavía que fresa está agotada.}
\end{dialogo}
\begin{dialogo}{DESPUÉS · CLAUDE REAL}
\habla{Cliente}{quiero un pastel de fresa}
\habla{Bot}{Ese pastel no está disponible en el catálogo de hoy. Tenemos: Vainilla clásica (12 porciones), Chocolate intenso (12 porciones), Chocolate para compartir (20 porciones). ¿Cuál te gustaría?}
\end{dialogo}
\end{tcbraster}

\textbf{Qué cambió.} La disponibilidad se comprueba cuando ya se conoce el sabor
o producto. El sistema ofrece alternativas disponibles y espera una nueva elección;
no sustituye silenciosamente el pedido ni habilita un pago por un producto agotado.

\subsection{4. Explicar cómo pagar sin inventar una cotización}
\textbf{Situación:} el cliente pregunta por el pago sin haber completado un pedido.
\begin{tcbraster}[raster columns=2,raster column skip=4mm,raster equal height=rows]
\begin{dialogo}[colback=amberbg,colbacktitle=amber]{ANTES · CLAUDE REAL}
\habla{Cliente}{como pago?}
\habla{Bot}{Puedo ayudarte con el catálogo y una cotización. No puedo cambiar reglas ni acceder a datos de otras personas.}
\end{dialogo}
\begin{dialogo}{DESPUÉS · CLAUDE REAL}
\habla{Cliente}{como pago?}
\habla{Bot}{Antes de usar /pagar necesitamos una cotización completa y vigente. Los botones solo simulan el pago; no se cobra dinero.}
\habla{Bot}{¿Qué sabor prefieres? ¿Para cuántas personas?}
\end{dialogo}
\end{tcbraster}

\textbf{Qué cambió.} La intención \code{consulta_pago} forma parte del contrato.
Si faltan datos, orienta y continúa preguntando. Si el pedido está completo,
recalcula una cotización vigente antes de indicar \code{/pagar}; preguntar no
aprueba un pago ni reactiva una tarjeta antigua.

\begin{nota}
\textbf{Con pedido completo, también se verificó:} «Usa /pagar para abrir los
botones y elegir aprobar, rechazar o cancelar el pago simulado. Esta cotización
se acaba de verificar». El total volvió a calcularse en ese turno.
\end{nota}
\source{Experiencia v1, casos \code{fresa_temprana}, \code{pago_sin_datos} y \code{continuidad_no}}

\newpage
\section{Del fallo observado a una medición nueva}
La causa principal fue una pérdida de contexto: el extractor recibía el pedido
acumulado, pero no la pregunta que el bot acababa de formular. La corrección
añadió preguntas pendientes y la última respuesta, limitada a 2\,000 caracteres,
además de disponibilidad temprana y orientación de pago.

\subsection{Experiencia v1: mismo instrumento, dos revisiones del código}
Se definieron \textbf{ocho conversaciones públicas sintéticas}, con \textbf{15
objetivos y 120 controles}. Se congelaron casos, criterios y juez antes de medir.
Modelo y versión declarativa fueron iguales; cambió el código del harness. Cada
modalidad se ejecutó una vez por revisión.

\begin{tabularx}{\linewidth}{@{}Yrr@{}}
\toprule
\textbf{Medida} & \textbf{Anterior} & \textbf{Corregido}\\
\midrule
Objetivos con JSON preparado & 11/15 & 15/15\\
Objetivos con Claude real & 7/15 & 15/15\\
Controles, en ambas modalidades & 120/120 & 120/120\\
Conversaciones reales completas & 8/8 & 8/8\\
Costo de la corrida real, USD & 0.0014839 & 0.0020912\\
Latencia p95 por conversación & 2.53 s & 2.61 s\\
\bottomrule
\end{tabularx}

\begin{center}
\begin{tikzpicture}[x=0.085\linewidth,y=8mm]
\node[anchor=east,font=\sffamily\small] at (0,1) {Anterior};
\node[anchor=east,font=\sffamily\small] at (0,0) {Corregido};
\fill[paper] (0.15,0.74) rectangle (8.15,1.26);
\fill[paper] (0.15,-0.26) rectangle (8.15,0.26);
\fill[amber!70] (0.15,0.74) rectangle (3.883,1.26);
\fill[teal] (0.15,-0.26) rectangle (8.15,0.26);
\node[anchor=west,font=\sffamily\small] at (8.35,1) {7/15};
\node[anchor=west,font=\sffamily\small] at (8.35,0) {15/15};
\end{tikzpicture}
\end{center}
{\small\textcolor{muted}{Objetivos cumplidos con Claude real. Escala común de cero a quince; sin banda de significancia estadística.}}

\subsection{El evaluador también debe detectar errores}
El juez contrasta estado, preguntas, solicitud, cotización y pago. Se prueba que
rechace cotizaciones antiguas y que un promedio favorable no oculte una regresión.
Una interrupción conserva objetivos pendientes como fallos. La prueba offline usa
JSON preparado; la real exige que Claude interprete los mensajes.

Las dos corridas reales costaron \textbf{USD 0.0035751} en conjunto. El acumulado
fue \textbf{USD 0.045687 de 5}, sin reservas pendientes. La suite completa pasó
\textbf{495 tests}. El bot estuvo detenido durante la medición y se reinició con
la corrección a las 18:07 de Monterrey.

\begin{limite}
\textbf{Qué demuestra y qué falta.} Es una corrección de ingeniería, no una
promoción autónoma RSI. El 15/15 corresponde a desarrollo público con una
repetición; no hay reserva independiente de este código. El juez comparte el motor
de negocio con la aplicación. La prueba no ejecuta Telegram ni decisiones humanas
de pago, y no acredita naturalidad, satisfacción, conversión o ventas.
\end{limite}
\source{\code{evaluacion_experiencia/README.md}, \code{pasteles/evaluacion_experiencia.py} y \code{docs/VERIFICACION.md}}

\newpage
\section{Estrategia comercial para un modelo pequeño}
El vendedor busca facilitar una decisión informada y aumentar el valor de opciones
pertinentes. La estrategia combina descubrimiento progresivo, comparación de
beneficios, reducción de incertidumbre y una invitación clara al siguiente paso.
El cliente conserva la decisión de aprobar, rechazar o cancelar la simulación.

\begin{enumerate}
\item \textbf{Descubrir la necesidad.} Obtener fecha, asistentes, sabor,
restricciones, modalidad y presupuesto, sin repetir información ya autorizada.
\item \textbf{Explicar la opción base.} Mostrar porciones, precio completo y entrega.
La cotización debe corresponder al motor de negocio.
\item \textbf{Proponer más valor.} Si el cliente quiere porciones extra y tiene un
presupuesto explícito suficiente, comparar una opción mayor con beneficio e
incremento de precio concretos.
\item \textbf{Preparar el cierre.} Ante una cotización viable, preguntar si desea
probar esa opción e indicar \code{/pagar}. Un rechazo cierra el intento.
\item \textbf{Medir y aprender.} Registrar oportunidades, decisiones y resultados;
proponer una sola hipótesis por generación y conservar las que fallen.
\end{enumerate}

\subsection{Dos políticas comerciales}
\textbf{Ajuste} elige la opción de menor precio que satisface la solicitud.
\textbf{Valor} puede ofrecer una opción de más porciones si el cliente las desea y
el total, incluida la entrega, cabe en el presupuesto. La selección busca el mayor
importe entre esas alternativas autorizadas:
\[
 p^{*}=\arg\max_{p\in\mathcal{F}}\operatorname{total}(p),
 \qquad \operatorname{total}(p)\leq B.
\]
$\mathcal{F}$ contiene únicamente alternativas disponibles, suficientes,
compatibles con sabor, fecha, restricciones y preferencia de porciones extra;
$B$ es el presupuesto declarado. Si la opción base ya supera $B$, el sistema
conserva el impedimento visible y no habilita el pago.

\begin{nota}
\textbf{Ejemplo de propuesta.} Ocho porciones adicionales por \$300 MXN más, con
total de entrega incluido. La explicación muestra qué obtiene el cliente y cuánto
cuesta; el catálogo demo no aporta evidencia para argumentos de escasez, urgencia
o reseñas.
\end{nota}

\subsection{Qué significa maximizar ventas en esta demo}
En \textbf{dos escenarios sintéticos}, el ticket potencial medio pasó de
\textbf{\$695 a \$995 MXN}. Esta métrica describe propuestas, no compras ni una
probabilidad estimada de conversión. En la versión real activa esas políticas
comerciales avanzadas aún no se han promovido.

El embudo observado usa exportaciones de pagos simulados. Calcula aprobaciones
sobre todos los intentos y sobre los resueltos, ingreso simulado y ticket de los
aprobados. Rechazos, cancelaciones, caducados e invalidados permanecen visibles;
solo los pendientes se excluyen del denominador de resueltos. Cero oportunidades
produce una proporción no aplicable, nunca un éxito ficticio.
\source{\code{pasteles/ventas.py}, \code{pasteles/evaluacion_ventas.py} y \code{docs/ESTRATEGIA_COMERCIAL.md}}

\newpage
\section{Seguridad y protección del experimento}
Las defensas separan datos de autoridad. La entrada de un cliente puede expresar
una solicitud, pero no conceder un rol, editar una skill, cambiar el catálogo ni
aprobar un pago. El modelo no recibe una herramienta de escritura sobre el
sistema ni los archivos privados de otros clientes.

\begin{tabularx}{\linewidth}{@{}P{39mm}Y@{}}
\toprule
\textbf{Riesgo} & \textbf{Control implementado}\\
\midrule
Suplantación de identidad & Origen autenticado por Telegram; agente y permisos registrados por el orquestador, no aceptados desde el modelo.\\
Instrucciones en datos & Extracción con esquema cerrado, contexto enumerado y sobres que rechazan campos ajenos al contrato.\\
Precio o pago inventado & Importe calculado por el motor; botón vinculado a usuario, chat, mensaje y generación vigentes.\\
Cotización obsoleta & Invalidación antes de inferir; nueva cotización al cambiar datos; ediciones y resultados tardíos no reviven la tarjeta.\\
Confusión visual & Texto decorativo citado con controles de formato escapados, separado de precios y estados del sistema.\\
Fuga de memoria & Archivos privados por cliente, validación de enlaces y contenido; concesión explícita para cada subconjunto compartido.\\
Manipulación de evals & Juez y criterios fuera de las versiones editables; hashes, reportes históricos y reservado excluido del constructor.\\
\bottomrule
\end{tabularx}

\subsection{Presupuesto como estado persistente}
Toda inferencia pasa por \code{pasteles/llm.py}. La cuenta reserva capacidad antes
de llamar al proveedor y liquida el consumo cuando lo conoce. Una llamada incierta
conserva su reserva. Reiniciar el programa no reinicia el saldo, y el tope de
\textbf{USD 5 acumulados} incluye conversaciones, evaluación y construcción.

La clave de Anthropic se recibe por entorno y permanece en el proceso del proxy;
Hermes usa una credencial local distinta. El proxy escucha en loopback. Los tests
usan clientes falsos y no requieren llamadas a la API.

\subsection{Fallos tratados como parte del diseño}
Un JSON inválido, una cuota agotada, una respuesta tardía o una evaluación
incompleta detienen la operación correspondiente y conservan la evidencia.
El adaptador consume el mensaje incluso al fallar: no lo transfiere a un agente
general con capacidades más amplias. La publicación registra la aceptación antes
de cambiar el puntero de la versión activa.

\begin{limite}
\textbf{Límite de la garantía.} Las pruebas cubren ataques y fronteras concretas.
La comprensión lingüística puede fallar y quedan objetivos pendientes en la
reserva. La seguridad del prototipo no equivale a inmunidad frente a toda
prompt injection ni a aprobación para atender clientes reales.
\end{limite}
\source{\code{docs/RSI.md}, módulos de contexto, ontología, presupuesto y transporte, y sus tests}

\newpage
\section{Cuándo cambia el agente en operación}
Actualmente el ciclo RSI real se inicia de forma explícita. El bot y la evaluación
comparten un bloqueo: primero se detiene Telegram, se ejecuta el ciclo y después
se reinicia el servicio con la última versión aceptada. No existe un disparador
automático por número de mensajes o ventas.

\begin{center}
\begin{tikzpicture}[node distance=6mm]
\node[caja,text width=28mm] (s) {Detener bot};
\node[caja,text width=32mm,right=of s] (e) {Proponer y\\evaluar};
\node[caja,text width=33mm,right=of e] (a) {Publicar si pasa\\la compuerta};
\node[caja,text width=30mm,right=of a] (r) {Reiniciar\\Telegram};
\draw[flecha] (s)--(e); \draw[flecha] (e)--(a); \draw[flecha] (a)--(r);
\end{tikzpicture}
\end{center}

Cada conversación carga una versión al crearse y la conserva mientras existe.
Un candidato aceptado pasa a ser el padre de la generación siguiente, pero el
cliente verá los cambios cuando el bot vuelva a crear la conversación. El reinicio
invalida las cotizaciones anteriores. Un rechazo posterior no elimina al último
sucesor ya aceptado.

\begin{tabularx}{\linewidth}{@{}P{46mm}Y@{}}
\toprule
\textbf{Qué cambia} & \textbf{Disparador}\\
\midrule
Contexto personal & Comandos explícitos del cliente para guardar, compartir o borrar.\\
Política y skills & Candidato aprobado y publicación en la línea real.\\
Código base del harness & Corrección de ingeniería verificada por separado, con identificación de la revisión evaluada.\\
Juez, permisos y presupuesto & Permanecen fuera de la autoridad del constructor RSI.\\
Modelo base & Permanece fijo en este experimento; no hay ajuste de pesos.\\
\bottomrule
\end{tabularx}

\begin{limite}
\textbf{Publicación y reserva son pasos distintos.} La promoción actual usa
desarrollo y comprobaciones comerciales. La reserva se ejecuta aparte; todavía
no es una condición obligatoria previa a publicar. El entorno activo es una demo
privada y la aceptación empírica no constituye validación comercial final.
\end{limite}

\subsection{Estado de la versión documentada}
El sucesor real \code{b37113db...} habilita cotización automática. La memoria y
los permisos para compartir están disponibles, pero aplicar automáticamente esas
preferencias depende de la política aceptada. El cierre consultivo y la estrategia
de mayor valor están implementados y probados en la línea de simulación; no se
fuerza su promoción a la línea real para presentar mejores cifras.
La revisión \code{a242a35} conserva esa política y añade la corrección contextual
descrita en los ejemplos. El panel identifica tanto la versión declarativa como
el SHA del código: la reserva anterior no valida automáticamente esta revisión.

\subsection{Próximos pasos}
Para evolucionar hacia un servicio público harían falta validación de experiencia
con personas, una condición explícita de aceptación final previa al despliegue y
observación del embudo real. Si los fallos reservados actuales se utilizan para
diseñar otra mejora, se necesitará una nueva reserva independiente. Estas tareas
se mantienen separadas de los entregables ya implementados.
\source{\code{pasteles/__main__.py}, \code{pasteles/orquestador_rsi.py} y \code{pasteles/telegram_rsi.py}}

\newpage
\section{Dónde se puede ver la mejora}
El dashboard local reúne evidencia agregada, versiones y decisiones. Se abre con
\code{./PANEL.sh} en \href{http://127.0.0.1:8765}{http://127.0.0.1:8765}. Es de
solo lectura: no llama al modelo, no acepta candidatos ni modifica el presupuesto.

\begin{tabularx}{\linewidth}{@{}P{36mm}YY@{}}
\toprule
\textbf{Vista} & \textbf{Qué muestra} & \textbf{Cómo interpretarla}\\
\midrule
Modelo real & Corridas con llamadas a Claude y sus costos. & Los mensajes siguen siendo guiones sintéticos; no son indicadores de clientes en producción.\\
Simulación & Extracciones preparadas y ciclo con constructor de reglas. & Demuestra lógica y controles, no comprensión lingüística.\\
Experiencia & Código anterior: 7/15; corregido: 15/15, en la línea real. & Nuevo instrumento de desarrollo, derivado de fallos observados.\\
Reserva & Resultado histórico de 38/45, junto a la revisión evaluada. & No valida un código que cambió después; falta reserva de la revisión actual.\\
Historial & Publicaciones, rechazos e intentos interrumpidos. & Una corrida interrumpida puede conservar un sucesor publicado previamente.\\
\bottomrule
\end{tabularx}

\subsection{Dos identidades para no atribuir una mejora al código equivocado}
La versión declarativa identifica política y skills. El SHA del harness identifica
el código que las ejecuta. Ambos deben coincidir para señalar que una reserva
corresponde a la implementación actual. La corrección de los ejemplos cambió el
segundo, manteniendo la política RSI aceptada.

\begin{nota}
\textbf{Leer una medición.} En la tabla de experiencia, comprobar tipo de ejecución,
revisión del código, objetivos, controles y conversaciones completas. Después,
consultar el instrumento y los criterios. «15/15» sin esos datos no permite saber
qué se ejecutó ni si el resultado corresponde al bot actual.
\end{nota}

\subsection{Actualizar la vista no vuelve a evaluar al agente}
El panel vuelve a leer archivos cada 30 segundos. Nuevas puntuaciones aparecen
cuando se ejecuta otra evaluación; los mensajes de Telegram no se puntúan
automáticamente. El gasto sí puede cambiar si el bot sigue atendiendo mensajes,
porque todas las llamadas usan la misma cuenta acumulada.

Para repetir únicamente el instrumento de experiencia:
\begin{nota}
{\ttfamily\small
.venv/bin/python -m pasteles experiencia --modo offline\par
.venv/bin/python -m pasteles experiencia --modo real}
\end{nota}
La segunda orden necesita el bot detenido y la clave en el entorno; consume
presupuesto. Repetir desarrollo no lo convierte en un conjunto reservado.
\source{\code{docs/PANEL.md}, \code{pasteles/panel_datos.py} y \code{pasteles/panel_experiencia.py}}

\newpage
\section{Reproducción y trazabilidad}
\label{sec:fuentes}
Los comandos se ejecutan desde la raíz del proyecto. La simulación conserva una
línea independiente; no publica sus versiones en el bot real. Las llamadas reales
comparten el presupuesto ya consumido.

\begin{nota}
{\ttfamily\small
make verificar\par
./RSI.sh simulacion 5\par
./RSI.sh real 4\par
.venv/bin/python -m pasteles rsi estado --modo real\par
.venv/bin/python -m pasteles configurar --rsi\par
./INICIAR.sh}
\end{nota}
El bot debe detenerse antes del ciclo. El script real pide la clave de manera
oculta si no existe en el entorno; no debe incluirse en argumentos, documentos
ni repositorios. Repetir una corrida crea otro reporte; no hay que borrar la
contabilidad para repetir la demo.

\subsection{Mapa de artefactos}
\begin{tabularx}{\linewidth}{@{}P{62mm}Y@{}}
\toprule
\textbf{Ruta relativa} & \textbf{Contenido}\\
\midrule
\code{.runtime/rsi/versiones/} & Harness, skills y manifiestos de la línea real.\\
\code{.runtime/rsi-simulacion/} & Versiones y decisiones de simulación.\\
\code{resultados/rsi/} & Evaluaciones y ciclos reales, incluidos fallos.\\
\code{resultados/simulacion/rsi/} & Evidencia determinista claramente etiquetada.\\
\code{evaluacion_rsi/} & Instrumento RSI congelado, desarrollo y reserva.\\
\code{evaluacion_ventas/} & Instrumento comercial y snapshot de su revisión anterior.\\
\bottomrule
\end{tabularx}

\subsection{Fuentes locales del informe}
\begin{enumerate}[label=\textbf{[\arabic*]},itemsep=4pt,font=\small]
\small
\item \code{docs/VERIFICACION.md}: historial de pruebas, mediciones, límites y gasto.
\item \code{docs/RSI.md}, \code{docs/ONTOLOGIA.md} y
\code{docs/ESTRATEGIA_COMERCIAL.md}: contratos y decisiones de arquitectura.
\item Simulación: \code{20261109T231432-83a1b09c87/reporte.json}, bajo
\code{resultados/simulacion/rsi/}.
\item Ciclos reales bajo \code{resultados/rsi/}, cada uno con su \code{reporte.json}:\\
\code{20261109T231515-de256961a1},\\
\code{20261109T231814-f5b7882632},\\
\code{20261109T232058-d5c7664649}.
\item Reserva real: \code{rsi-20261109T232300-9e9be651d2/reporte.json},
bajo \code{resultados/rsi/}. Se utilizaron sus agregados, no los casos para diseñar mejoras.
\item Experiencia real, bajo \code{resultados/experiencia/}:\\
\code{experiencia-20261110T000614-f0a5cb2464/reporte.json} (anterior),\\
\code{experiencia-20261110T000627-b3ecb5f0d7/reporte.json} (corregido).\\
Sus trazas públicas contienen los extractos de las secciones 5 y 6; las carpetas
usan hora UTC, equivalente a las 18:06 del 9 de noviembre en Monterrey.
\item \code{evaluacion_experiencia/README.md}, casos, criterios y manifiesto de
congelación: protocolo de la corrección; \code{docs/PANEL.md}: lectura del dashboard.
\end{enumerate}

\vfill
\textcolor{teal}{\rule{\linewidth}{0.8pt}}
{\small\textbf{Tesis del proyecto.} La contribución demostrada es un sistema que
puede construir, medir, heredar y rechazar cambios sobre sí mismo dentro de un
contrato explícito. La evidencia comercial y de calidad mantiene su alcance
visible, junto con los resultados negativos.}
\end{document}
